% !TEX root = 04.tex
\documentclass[a4paper,10pt]{article}

\usepackage{pdfsync}
\usepackage[utf8]{inputenc}
\usepackage{microtype}
\usepackage{mathtools}
\usepackage{amsthm,amsfonts,amsmath,amssymb}
\usepackage{mathabx}
% \usepackage{MnSymbol}
\usepackage{hyperref}
\usepackage{bbold}
\usepackage[textwidth=2cm,textsize=small]{todonotes}
\usepackage{subcaption}
\usepackage{enumerate}
\usepackage{xspace}
\usepackage{stmaryrd}
\usepackage{verbatim}
\usepackage{cleveref}

\hypersetup{
    colorlinks,
    linkcolor={red!50!black},
    citecolor={blue!50!black},
    urlcolor={blue!80!black}
}

\author{}

\newcommand{\ignore}[1]{}

\newcommand{\st}{s.t.}
\newcommand{\ie}{i.e.}
\newcommand{\eg}{e.g.}

\newcommand{\sep}{\;|\;}
\renewcommand{\rule}[2]{\dfrac{#1}{#2}}
\newcommand{\set}[1]{\left\{#1\right\}}
\newcommand{\setof}[2]{\set{#1 \;\middle|\; #2}}
\newcommand{\tuple}[1]{\left( #1 \right)}
\newcommand{\pair}[2]{\tuple {#1, #2}}
\newcommand{\closure}[3]{\tuple {#1, #2, #3}}
\newcommand{\sem}[1]{\llbracket{#1}\rrbracket}
\newcommand{\powerset}[1]{\mathcal P(#1)}
\newcommand{\map}[2]{\mathsf{map}\; #1\; #2}
\newcommand{\filter}[2]{\mathsf{filter}\; #1\; #2}
\newcommand{\floor}[1]{\left\lfloor #1 \right\rfloor}
\newcommand{\PreRel}[2]{\mathsf{Pre}_{#1}(#2)}
\newcommand{\Pre}[1]{\PreRel {} {#1}}
\newcommand{\PreStar}[1]{\mathsf{Pre}^*(#1)}
\newcommand{\Post}[1]{\mathsf{Post}(#1)}
\newcommand{\PostStar}[1]{\mathsf{Post^*}(#1)}
\newcommand{\Conf}{\mathsf{Conf}}
\newcommand{\upwardclosure}[1]{#1\uparrow}
\newcommand{\card}[1]{\left| #1 \right|}
\newcommand{\lang}[1]{L(#1)}
\newcommand{\goesto}[1]{\xrightarrow{#1}}
\newcommand{\e}{\varepsilon}

\newcommand{\N}{\mathbb N}
\newcommand{\Z}{\mathbb Z}
\newcommand{\Q}{\mathbb Q}

% LTL
\newcommand{\true}{\mathbf{true}}
\newcommand{\false}{\mathbf{false}}
\newcommand{\X}{\mathbf X}
\newcommand{\U}{\mathbf U}
\newcommand{\Release}{\mathbf R}
\newcommand{\F}{\mathbf F}
\newcommand{\G}{\mathbf G}

% CTL
\newcommand{\EX}{\mathbf{EX}}
\newcommand{\AX}{\mathbf{AX}}
\newcommand{\EF}{\mathbf{EF}}
\newcommand{\AF}{\mathbf{AF}}
\newcommand{\EG}{\mathbf{EG}}
\newcommand{\AG}{\mathbf{AG}}
\newcommand{\EU}[2]{\mathbf E(#1 \U #2)}
\newcommand{\AU}[2]{\mathbf A(#1 \U #2)}
\newcommand{\ER}[2]{\mathbf E(#1 \Release #2)}
\newcommand{\AR}[2]{\mathbf A(#1 \Release #2)}

\newcommand{\NL}{\textsf{NL}}
\newcommand{\PTIME}{\textsf{PTIME}}
\newcommand{\PSPACE}{\textsf{PSPACE}}

\makeatletter

\def\dasharrowfill@#1#2#3#4{%
        $\m@th
        \thickmuskip0mu
        \medmuskip\thickmuskip
        \thinmuskip\thickmuskip
        \relax
        #4#1\mkern2mu
        \xleaders\hbox{$#4\mkern2mu#2\mkern2mu$}\hfill
        \mkern2mu
        #3$%
}

\def\dashleftarrowfill@{\dasharrowfill@\leftarrow\relbar\relbar}
\def\dashrightarrowfill@{\dasharrowfill@\relbar\relbar\rightarrow}
\def\dashleftrightarrowfill@{\dasharrowfill@\leftarrow\relbar\rightarrow}
\def\dashLeftarrowfill@{\dasharrowfill@\Leftarrow\Relbar\Relbar}
\def\dashRightarrowfill@{\dasharrowfill@\Relbar\Relbar\Rightarrow}
\def\dashLeftrightarrowfill@{\dasharrowfill@\Leftarrow\Relbar\Rightarrow}

\providecommand*\xdashleftarrow[2][]{%
  \ext@arrow 0055{\dashleftarrowfill@}{#1}{#2}}
\providecommand*\xdashrightarrow[2][]{%
  \ext@arrow 0055{\dashrightarrowfill@}{#1}{#2}}
\providecommand*\xdashleftrightarrow[2][]{%
  \ext@arrow 0055{\dashleftrightarrowfill@}{#1}{#2}}
\providecommand*\xdashLeftarrow[2][]{%
  \ext@arrow 0055{\dashLeftarrowfill@}{#1}{#2}}
\providecommand*\xdashRightarrow[2][]{%
  \ext@arrow 0055{\dashRightarrowfill@}{#1}{#2}}
\providecommand*\xdashLeftrightarrow[2][]{%
  \ext@arrow 0055{\dashLeftrightarrowfill@}{#1}{#2}}

\def\slashedarrowfill@#1#2#3#4#5{%
$\m@th\thickmuskip0mu\medmuskip\thickmuskip\thinmuskip\thickmuskip
  \relax#5#1\mkern-7mu%
  \cleaders\hbox{$#5\mkern-2mu#2\mkern-2mu$}\hfill
  \mathclap{#3}\mathclap{#2}%
  \cleaders\hbox{$#5\mkern-2mu#2\mkern-2mu$}\hfill
  \mkern-7mu#4$%
}

\def\rightslashedarrowfill@{%
  \slashedarrowfill@\relbar\relbar{\raisebox{1.2pt}{$\scriptscriptstyle\diagup$}}\rightarrow}
\newcommand\xslashedrightarrow[2][]{%
  \ext@arrow 0055{\rightslashedarrowfill@}{#1}{#2}}
\makeatother

\newtheorem{exercise}{Exercise}

\newenvironment{solution}{\begin{proof}[\protect{\textnormal{\emph{Solution.}}}]}{\end{proof}}

\let\oldqedhere\qedhere

\newif\ifstartedinmathmode
\renewcommand*{\qedhere}{%
  \relax\ifmmode\startedinmathmodetrue\else\startedinmathmodefalse\fi%
  \ifstartedinmathmode\tag*{\protect\oldqedhere}\else\ifinlist{\hfill\oldqedhere}{\oldqedhere}\fi%
}

\crefname{equation}{}{}
\crefname{exercise}{Exercise}{Exercises}


\title{Formal verification - Tutorial 04}
\date{16/03/2026}

\begin{document}

\maketitle

\newcommand{\lossyto}{\Rightarrow}
\section*{Communicating finite-state machines}

Let $S$ be CFSM with (perfect) transition relation $\to$.

\begin{exercise}
	Show that the reachability problem for CFSMs is undecidable.
\end{exercise}

\begin{solution}
	We can simulate the tape of a Turing machine using a single channel.
	A special symbol is used to mark the position of the head.
	The head movement are simulated by rotations of the channel contents.
\end{solution}

We consider a model of CFSM comprising a finite set of processes $P$ (each with a finite set of local states)
and a finite set of channels $C$.
%
Each channel is a FIFO queue of messages connecting two processes (point-to-point communication).
%
The \emph{communication topology} of a CFSM is an undirected graph
where vertices are processes and (undirected) edges are channels.

\begin{exercise}
	Show that the reachability problem for CFSMs with is decidable \emph{over tree communication topologies}.
\end{exercise}

\begin{solution}
	When the communication topology is a tree,
	bounded channels suffice to reach all reachable configurations.
	%
	The idea is that we can always keep the channels to size zero or one,
	by always executing a transmision immediately followed by the corresponding reception.
\end{solution}

\section*{Well-quasi orders}

A binary relation $\preceq$ on a set $X$ 
is \textbf{quasi-order}
if it is reflexive and transitive.
An antichain is a subset of $X$ 
whose elements are pairwise incomparable w.r.t.~$\preceq$.

\begin{exercise}
    Let $(X, \preceq)$ be a quasi-order. Prove that the following conditions
    are equivalent.
    \begin{enumerate}%[label=($\alph*$)]
        \item[a)] Any infinite sequence $a_1, a_2, \dots$ of elements of $X$
        contains a \emph{domination} $a_i \preceq a_j$
        for some indexes $i < j$.

        \item[b)] $X$ does not contain any infinite antichain and is \emph{well-founded}
        (i.e., there are no infinite strictly decreasing sequences).

        \item[c)] Any infinite sequence $a_1, a_2, \dots$ of elements of $X$
        contains an infinite nondecreasing subsequence, i.e., there are indices
        $i_1 < i_2 < \ldots$ with $a_{i_1} \preceq a_{i_2} \preceq \ldots$.

        \item[d)] Every upward closed set is the upward closure of a finite set.
        (The \emph{upward closure} of a set $X$
        is the set of elements that dominate some element of $X$;
        $X$ is \emph{upward closed} if it is equal to its upward closure.)

        \item[e)] Every nondecreasing chain of upward closed sets $U_1 \subseteq U_2 \subseteq \cdots \subseteq X$ is finite.

    \end{enumerate}
\end{exercise}
\noindent
Whenever any of the conditions above holds we call 
the relation $\preceq$ a \textbf{well quasi-order} or \textbf{WQO}.
% \begin{solution}
    % \noindent
    % ($a \Rightarrow b$) Well-foundedness is obvious.
    % Now consider an infinite antichain $A \subseteq X$
    % and arrange its elements into an infinite sequence.
    % %
    % By the assumption we can find a domination, contradicting that $A$ is an antichain.

    % \noindent
    % ($b \Rightarrow a$) Consider the set of minimal elements of the infinite sequence.
    % This set is nonempty by well-foundedness and it is finite since it is an antichain.
    % Take any element in the sequence not from this finite set.
    % It is preceded by a smaller element in the set, yielding a domination.

    % ($a \Rightarrow c$) By way of contradiction,
    % assume each $x_i$ is the beginning of chains of finite length only.
    % %
    % From $x_0$, select one such maximal chain, ending at some $x_{i_0}$.
    % From $x_{i_0+1}$, select one such maximal chain, ending at some $x_{i_1}$, etc.
    % %    
    % We have constructed an infinite subsequence of elements $x_{i_0}, x_{i_1}, \dots$
    % which are not dominated by \emph{any} element to their right,
    % leading to a contradiction.

    % \noindent
    % ($c \Rightarrow a$) This is trivial.

    % \noindent
    % ($b \Rightarrow d$) Let $U \subseteq X$ be upward closed and consider the subset of its elements $M$
    % which are not dominated by any element in $U$ (minimal elements).
    % Since there order is well-founded, $U$ is the upward closure of $M$.
    % The latter set is an antichain, thus finite by assumption.

    % \noindent
    % ($d \Rightarrow b$) Clearly antichains must be finite.
    % Regarding well-foundedness, by way of contradiction assume we have
    % an infinite strictly decreasing sequence $x_1 \succ x_2 \succ \cdots$.
    % %
    % Then the upward closure of its elements is not the upward closure of any finite set,
    % which is a contradiction.

    % \noindent
    % ($d \Rightarrow e$) Consider the upward closed set $U_1 \cup U_2 \cup \cdots$.
    % By assumption, it is the upward closure of finitely many elements $B = \set{x_1, \dots, x_m}$.
    % Take the first $U_n$ \st~$B \subseteq U_n$ and we have $U_n = U_{n+1} = \cdots$.

    % \noindent
    % ($e \Rightarrow d$) Let $U \subseteq X$ be an upward closed set and consider the set of its minimal elements
    % $M = \set{x_1, x_2, \dots}$ (if there are ties just select arbitrarily an element from each minimal equivalence class).
    % Now consider the chain generated by the upward closure of $\set{x_1}$, $\set{x_1, x_2}$, etc.
    % By assumption, this chain is finite, therefore there is $n$ \st~the upward closure of $\set{x_1, \dots, x_n}$
    % is the same as that of $M$ itself, which is just $X$.
% \end{solution}

\begin{exercise}
	Show that every upward-closed language of words (w.r.t.~the subword ordering) is regular.
\end{exercise}

% \begin{solution}
% \end{solution}

\section*{Lossy channel systems}

Let $S$ be a CFSM with (perfect) transition relation $\to$
and lossy transition relation $\lossyto$.
%
Recall that $c \lossyto d$ if there are configurations $c', d' \in \Conf$ such that
$c' \sqsubseteq c$, $c' \to d'$, and $d \sqsubseteq d'$.
%
Consider the following predecessor operators:
%
\begin{align*}
	\PreRel \to T
		&:= \setof {c \in \Conf}{\exists c' \in T, c \to c'}, \\
	\Pre T
		&:= \setof {c \in \Conf}{\exists c' \in T, c \Rightarrow c'},
\end{align*}
%
Notice that $\Pre T$ can be infinite even when $T$ only contains a single configuration.
%
\begin{exercise}
	Let $T$ be a regular set of configurations.
	Show that the following sets are effectively regular:
	\begin{itemize}
		\item $\PreRel \to T$,
		\item $\upwardclosure T$,
		\item $\Pre T$,
		\item $\PreStar T$.
	\end{itemize}
\end{exercise}

\begin{solution}
	The set $\PreRel \to T$ can be obtained from $T$ by regularity-preserving operations
	such as finite unions, concatenations, and right language quotients.
	%
	From a finite automaton recognising $T$,
	we can construct a finite automaton recognising $\upwardclosure T$ by adding
	self-loops on all states for all symbols.
	%
	Straight from the definitions, we have
	\begin{align*}
		\Pre T = \upwardclosure{\PreRel \to {\upwardclosure T}}. \qedhere
	\end{align*}
\end{solution}

\begin{exercise}
	Show that the forward reachability set for lossy channel systems is regular.
	Can one \emph{compute} a finite automaton recognising it?
\end{exercise}


\bibliographystyle{plainurl}
\bibliography{bibliography}

\end{document}